\documentclass[11pt,a4paper]{article}
\usepackage[utf8]{inputenc}
\usepackage[T1]{fontenc}
\usepackage{lmodern}
\usepackage{amsmath}
\usepackage[margin=2cm]{geometry}
\usepackage{enumitem}
\pagestyle{plain}
\setlength{\parindent}{0pt}
\setlength{\parskip}{0pt}
\begin{document}
{\Large\bfseries Práctica de química}
\vspace{12pt}
\begin{enumerate}[label=\textbf{\arabic*.},leftmargin=*,itemsep=16pt,topsep=0pt]
\item \begin{minipage}[t]{\linewidth}
Una muestra tiene una masa de 650 mg. Expresa esa masa en gramos.

\vspace{7pt}
\begin{tabular}{@{}p{0.48\linewidth}p{0.48\linewidth}@{}}
a) 6.5 g & b) 0.65 g \\[5pt]
c) 65 g & d) 0.065 g \\
\end{tabular}
\end{minipage}

\item \begin{minipage}[t]{\linewidth}
Una lámina tiene un grosor de 0.12 mm. ¿Cuántos micrómetros (\(\mu\)m) mide?

\vspace{7pt}
\begin{tabular}{@{}p{0.48\linewidth}p{0.48\linewidth}@{}}
a) 12 \(\mu\)m & b) 0.00012 \(\mu\)m \\[5pt]
c) 1200 \(\mu\)m & d) 120 \(\mu\)m \\
\end{tabular}
\end{minipage}

\item \begin{minipage}[t]{\linewidth}
Una muestra contiene 1250 \(\mu\)g de una sustancia. Expresa su masa en mg.

\vspace{7pt}
\begin{tabular}{@{}p{0.48\linewidth}p{0.48\linewidth}@{}}
a) 0.125 mg & b) 125 mg \\[5pt]
c) 1.25 mg & d) 1250000 mg \\
\end{tabular}
\end{minipage}

\item \begin{minipage}[t]{\linewidth}
Un recipiente tiene 0.48 L de agua. Expresa el volumen en cm\textsuperscript{3}.

\vspace{7pt}
\begin{tabular}{@{}p{0.48\linewidth}p{0.48\linewidth}@{}}
a) 480 cm\textsuperscript{3} & b) 48 cm\textsuperscript{3} \\[5pt]
c) 0.48 cm\textsuperscript{3} & d) 4800 cm\textsuperscript{3} \\
\end{tabular}
\end{minipage}

\item \begin{minipage}[t]{\linewidth}
Una muestra está a \(-\)25 \({}^{\circ}\mathrm{C}\). Expresa su temperatura en K.

\vspace{7pt}
\begin{tabular}{@{}p{0.48\linewidth}p{0.48\linewidth}@{}}
a) 298.15 K & b) \(-\)248.15 K \\[5pt]
c) 25 K & d) 248.15 K \\
\end{tabular}
\end{minipage}


\newpage
\item \begin{minipage}[t]{\linewidth}
Una muestra está a 104 \({}^{\circ}\mathrm{F}\). Su temperatura disminuye 10 \({}^{\circ}\mathrm{C}\). Calcula la temperatura final en \({}^{\circ}\mathrm{F}\) y en K.

\vspace{7pt}
\begin{tabular}{@{}p{0.48\linewidth}p{0.48\linewidth}@{}}
a) 94 \({}^{\circ}\mathrm{F}\); 307.59 K & b) 86 \({}^{\circ}\mathrm{F}\); 303.15 K \\[5pt]
c) 86 \({}^{\circ}\mathrm{F}\); 243.15 K & d) 122 \({}^{\circ}\mathrm{F}\); 323.15 K \\
\end{tabular}
\end{minipage}

\item \begin{minipage}[t]{\linewidth}
Una pieza tiene una masa de 0.216 kg y ocupa 80 cm\textsuperscript{3}. Calcula su densidad en g/cm\textsuperscript{3} y en kg/m\textsuperscript{3}.

\vspace{7pt}
\begin{tabular}{@{}p{0.48\linewidth}p{0.48\linewidth}@{}}
a) 0.0027 g/cm\textsuperscript{3}; 2.7 kg/m\textsuperscript{3} & b) 27 g/cm\textsuperscript{3}; 27000 kg/m\textsuperscript{3} \\[5pt]
c) 2.7 g/cm\textsuperscript{3}; 2700 kg/m\textsuperscript{3} & d) 2.7 g/cm\textsuperscript{3}; 2.7 kg/m\textsuperscript{3} \\
\end{tabular}
\end{minipage}

\item \begin{minipage}[t]{\linewidth}
Se necesitan 0.24 L de un líquido de densidad 1.25 g/mL. ¿Qué masa tendrá en kg?

\vspace{7pt}
\begin{tabular}{@{}p{0.48\linewidth}p{0.48\linewidth}@{}}
a) 3 kg & b) 0.30 kg \\[5pt]
c) 0.0003 kg & d) 300 kg \\
\end{tabular}
\end{minipage}

\item \begin{minipage}[t]{\linewidth}
Tienes 0.084 kg de un líquido de densidad 1.20 g/mL. ¿Qué volumen ocupa en L?

\vspace{7pt}
\begin{tabular}{@{}p{0.48\linewidth}p{0.48\linewidth}@{}}
a) 0.1008 L & b) 0.70 L \\[5pt]
c) 0.070 L & d) 70 L \\
\end{tabular}
\end{minipage}

\item \begin{minipage}[t]{\linewidth}
Una pieza de 0.162 kg hace subir el agua de una probeta de 25 a 85 mL. Está completamente sumergida y no se disuelve. Calcula su densidad en g/cm\textsuperscript{3}.

\vspace{7pt}
\begin{tabular}{@{}p{0.48\linewidth}p{0.48\linewidth}@{}}
a) 2.7 g/cm\textsuperscript{3} & b) 1.91 g/cm\textsuperscript{3} \\[5pt]
c) 0.0027 g/cm\textsuperscript{3} & d) 6.48 g/cm\textsuperscript{3} \\
\end{tabular}
\end{minipage}


\newpage
\item \begin{minipage}[t]{\linewidth}
Un bloque macizo mide 30 mm \(\times\) 4 cm \(\times\) 0.02 m y tiene una masa de 0.0648 kg. Calcula su volumen en cm\textsuperscript{3} y su densidad en g/cm\textsuperscript{3}.

\vspace{7pt}
\begin{tabular}{@{}p{0.48\linewidth}p{0.48\linewidth}@{}}
a) 24 cm\textsuperscript{3}; 2.7 g/cm\textsuperscript{3} & b) 240 cm\textsuperscript{3}; 0.27 g/cm\textsuperscript{3} \\[5pt]
c) 2.4 cm\textsuperscript{3}; 27 g/cm\textsuperscript{3} & d) 24 cm\textsuperscript{3}; 0.0027 g/cm\textsuperscript{3} \\
\end{tabular}
\end{minipage}

\item \begin{minipage}[t]{\linewidth}
Un frasco vacío tiene una masa de 60 g. Lleno de un líquido de densidad 1.08 g/mL, su masa total es 0.384 kg. Se vacía, se seca y se llena al mismo nivel con otro líquido de 0.92 g/mL. Calcula el volumen en L y la nueva masa total en g.

\vspace{7pt}
\begin{tabular}{@{}p{0.48\linewidth}p{0.48\linewidth}@{}}
a) 0.30 L; 276 g & b) 0.03 L; 87.6 g \\[5pt]
c) 0.30 L; 336 g & d) 0.3556 L; 387.1 g \\
\end{tabular}
\end{minipage}

\item \begin{minipage}[t]{\linewidth}
Un bloque rectangular macizo tiene una masa de 0.432 kg y densidad de 2.70 g/cm\textsuperscript{3}. Su base mide 40 mm \(\times\) 0.05 m. Calcula la altura en mm.

\vspace{7pt}
\begin{tabular}{@{}p{0.48\linewidth}p{0.48\linewidth}@{}}
a) 0.8 mm & b) 8 mm \\[5pt]
c) 800 mm & d) 80 mm \\
\end{tabular}
\end{minipage}

\item \begin{minipage}[t]{\linewidth}
Un recipiente de 1.20 L contiene 450 cm\textsuperscript{3} de un líquido de densidad 0.80 g/mL. ¿Qué masa adicional en kg necesitas para llenarlo?

\vspace{7pt}
\begin{tabular}{@{}p{0.48\linewidth}p{0.48\linewidth}@{}}
a) 0.96 kg & b) 0.60 kg \\[5pt]
c) 0.36 kg & d) 600 kg \\
\end{tabular}
\end{minipage}

\item \begin{minipage}[t]{\linewidth}
Una muestra pasa de 278.15 K a 86 \({}^{\circ}\mathrm{F}\). Calcula el aumento de temperatura en \({}^{\circ}\mathrm{C}\) y la temperatura final en K.

\vspace{7pt}
\begin{tabular}{@{}p{0.48\linewidth}p{0.48\linewidth}@{}}
a) 25 \({}^{\circ}\mathrm{C}\); 298.15 K & b) 35 \({}^{\circ}\mathrm{C}\); 308.15 K \\[5pt]
c) 25 \({}^{\circ}\mathrm{C}\); 303.15 K & d) 81 \({}^{\circ}\mathrm{C}\); 354.15 K \\
\end{tabular}
\end{minipage}

\end{enumerate}
\end{document}
